\chapter{Technical Appendix}

\section{Main Pipeline Route}

The following condensed route can be used to inspect logs or explain the implementation during a demonstration:

\begin{enumerate}
  \item Check required programs and source all modules.
  \item Initialise output directories and the working database.
  \item Select single-resource or batch mode.
  \item Find or create the repository and create a new run.
  \item Detect the platform and normalize supplied notebook paths.
  \item Validate through Git or the Zenodo Records API.
  \item Fetch and store normalized metadata.
  \item Clone, pull, or download and extract the resource.
  \item Register supplied notebooks or discover \texttt{.ipynb} files.
  \item Check total and Python notebook counts.
  \item Build requirements and create a pyenv environment.
  \item Execute notebooks and classify errors.
  \item Compare original and executed outputs.
  \item Clean the environment and finalise the run.
\end{enumerate}

\section{Representative Database Queries}

\begin{lstlisting}[language=SQL,caption={Inspect normalized metadata by platform.}]
SELECT r.platform,
       r.repository,
       m.title,
       m.authors,
       m.license,
       m.doi
FROM repositories AS r
LEFT JOIN repository_metadata AS m
  ON m.repo_id = r.id
ORDER BY r.platform, r.repository;
\end{lstlisting}

\begin{lstlisting}[language=SQL,caption={Inspect repository run outcomes.}]
SELECT r.platform,
       rr.run_status,
       COUNT(*) AS runs
FROM repository_runs AS rr
JOIN repositories AS r
  ON r.id = rr.repository_id
GROUP BY r.platform, rr.run_status
ORDER BY r.platform, rr.run_status;
\end{lstlisting}

\section{Representative SPARQL Queries}

\begin{lstlisting}[caption={List resources and platforms.}]
PREFIX nf: <https://w3id.org/notebookfair#>
PREFIX rdfs: <http://www.w3.org/2000/01/rdf-schema#>

SELECT ?resource ?resourceLabel ?platformLabel
WHERE {
  ?resource a nf:RepositoryResource ;
            rdfs:label ?resourceLabel ;
            nf:hostedOnPlatform ?platform .
  ?platform rdfs:label ?platformLabel .
}
ORDER BY ?platformLabel ?resourceLabel
\end{lstlisting}

\begin{lstlisting}[caption={Constraint check for Zenodo typing.}]
PREFIX nf: <https://w3id.org/notebookfair#>

SELECT ?resource
WHERE {
  ?resource nf:hostedOnPlatform
      <https://w3id.org/notebookfair/platform/zenodo> ;
      a nf:GitRepositoryResource .
}
# Expected result: zero rows.
\end{lstlisting}

\begin{lstlisting}[caption={Trace an execution to a result.}]
PREFIX nf: <https://w3id.org/notebookfair#>

SELECT ?repository ?run ?execution ?result ?score
WHERE {
  ?repository nf:hasRun ?run .
  ?run nf:hasExecution ?execution .
  ?execution nf:hasReproducibilityResult ?result .
  ?result nf:reproducibilityScore ?score .
}
\end{lstlisting}

\section{Build Commands}

\begin{lstlisting}[language=bash,caption={Metadata test in the local container image.}]
docker build -f binder/Dockerfile -t notebookfair-test .
docker run --rm notebookfair-test bash tests/test_metadata.sh
\end{lstlisting}

\begin{lstlisting}[language=bash,caption={Build the pipeline knowledge graph.}]
python kg/fairjupyter/build_pipeline_kg.py
\end{lstlisting}

The graph builder writes the combined graph and build report to:
\begin{center}
\path{output/kg/notebookfair-pipeline.nt}\\
\path{output/kg/build-summary.json}
\end{center}

\section{Draft Completion Checklist}

\begin{itemize}
  \item[\texttt{[ ]}] Review final title and student administration details.
  \item[\texttt{[ ]}] Complete ontology domain, range, inverse, and datatype review.
  \item[\texttt{[ ]}] Freeze and run the cross-platform experiment sample.
  \item[\texttt{[ ]}] Add comparative tables and plots with denominators.
  \item[\texttt{[x]}] Implement hybrid notebook classification and human-review warnings.
  \item[\texttt{[ ]}] Evaluate notebook classification against a human-labelled sample.
  \item[\texttt{[ ]}] Replace interim evaluation statements with final results.
  \item[\texttt{[ ]}] Review references, language, and university formatting requirements.
  \item[\texttt{[ ]}] Obtain supervisor feedback and complete the declaration page.
\end{itemize}
